\section{Descripción de la Problemática}
 
\subsection{Situación actual}
 
Actualmente, la gestión de los procedimientos administrativos establecidos en el TUPA de la UNSAAC presenta diversas limitaciones. El sistema PLADDES fue desarrollado con tecnologías de generación anterior: \textbf{AdminLTE 2 + Bootstrap 3 + jQuery}, frameworks que ya no reciben actualizaciones activas de seguridad ni mejoras de rendimiento.
 
Muchos procesos aún dependen de mecanismos manuales, documentación física y procedimientos descentralizados que dificultan la eficiencia institucional. Los estudiantes, egresados y demás administrados enfrentan constantemente problemas como largas colas, desconocimiento de requisitos, retrasos en la validación documentaria, falta de seguimiento en tiempo real de sus expedientes y poca transparencia respecto al estado de sus trámites.
 
\subsection{Formulación del problema}
 
\textbf{Problema general:} ¿De qué manera una plataforma web transaccional puede optimizar la gestión de los procedimientos administrativos del TUPA de la UNSAAC, mejorando la eficiencia, transparencia y seguimiento de los trámites universitarios?
 
\textbf{Problemas específicos:}
\begin{itemize}
    \item ¿Cómo reducir el tiempo de atención y procesamiento de los trámites administrativos?
    \item ¿Cómo centralizar la información relacionada con requisitos, costos, plazos y estados de los procedimientos del TUPA?
    \item ¿Cómo mejorar el seguimiento y la trazabilidad de los expedientes en tiempo real?
    \item ¿Cómo optimizar la gestión documental y el flujo de trabajo entre las diferentes oficinas?
    \item ¿Cómo implementar mecanismos digitales que disminuyan errores manuales y mejoren la transparencia administrativa?
\end{itemize}
 
\subsection{Problemas identificados en PLADDES}
 
\begin{longtable}{>{\bfseries}p{1.2cm} p{12cm}}
\toprule
\textbf{Cód.} & \textbf{Problema} \\
\midrule
\endhead
P1 & \textbf{Deuda tecnológica acumulada:} el frontend está construido sobre AdminLTE 2 con Bootstrap 3 y jQuery, tecnologías obsoletas desde 2019. Esto dificulta el mantenimiento, la incorporación de nuevas funcionalidades y la corrección de vulnerabilidades de seguridad. \\[0.3cm]
P2 & \textbf{Experiencia de usuario deficiente en el formulario de solicitud:} el proceso de registro se presenta como un wizard de 4 pasos con una interfaz poco intuitiva, sin validación en tiempo real ni mensajes de error claros. \\[0.3cm]
P3 & \textbf{Seguimiento de trámites limitado:} la opción ``Verifique AQUÍ su trámite'' no ofrece vista de línea de tiempo, historial de estados ni información sobre la dependencia responsable en cada momento. \\[0.3cm]
P4 & \textbf{Verificación de pagos manual y sin integración:} el sistema exige adjuntar el voucher de pago, pero no existe verificación automática con los sistemas de caja de la UNSAAC. \\[0.3cm]
P5 & \textbf{Ausencia de notificaciones activas:} el único mecanismo de comunicación es el correo electrónico ingresado manualmente, sin confirmaciones automáticas al registrar ni alertas al cambiar el estado del trámite. \\[0.3cm]
P6 & \textbf{Problemas de disponibilidad:} el sistema muestra avisos del tipo ``NO REALICE NINGÚN TRÁMITE, ESTAMOS EN MANTENIMIENTO'' directamente en la página principal, generando desconfianza. \\[0.3cm]
P7 & \textbf{Catálogo de trámites no integrado:} los requisitos, costos y plazos del TUPA están disponibles únicamente como documento PDF externo, desvinculado del sistema. \\[0.3cm]
P8 & \textbf{Gestión documental sin estándar claro:} el sistema acepta múltiples formatos pero no guía al usuario sobre qué documentos específicos debe adjuntar según el tipo de trámite seleccionado. \\[0.3cm]
P9 & \textbf{Sin panel administrativo de gestión:} las dependencias no cuentan con un panel dedicado para monitorear su bandeja de expedientes pendientes, asignar responsables ni generar reportes de productividad. \\
\bottomrule
\end{longtable}
 
\subsection{Impacto del problema}
 
\begin{itemize}
    \item \textbf{Retrasos en la atención:} trámites observados por documentación incompleta que pudo prevenirse con guía contextual en el formulario.
    \item \textbf{Sobrecarga del personal} de la Unidad de Trámite Documentario por validaciones manuales de pagos y consultas telefónicas de seguimiento.
    \item \textbf{Insatisfacción de la comunidad universitaria,} especialmente de estudiantes de las filiales que dependen exclusivamente del canal virtual.
    \item \textbf{Riesgo de incumplimiento del silencio administrativo positivo} establecido en el TUPA (plazos de 5 a 30 días hábiles según el procedimiento).
    \item \textbf{Imposibilidad de auditar} el desempeño administrativo por ausencia de estadísticas y reportes automatizados.
\end{itemize}